%% I. Introduction: background, problem, gap, research questions, contributions, organisation. Input by access.tex.
\section{Introduction}
\label{sec:intro}
\PARstart{M}{editation} research has moved from mean differences between meditators and controls towards models of the dynamics of a session: the alternation between focused attention, mind wandering and meta-awareness \cite{Hasenkamp2012,Lutz2008,Christoff2016}, the regulation of the default mode network (DMN) \cite{Raichle2001,Brewer2011,Garrison2015}, and computational accounts in terms of predictive processing \cite{Laukkonen2021,SandvedSmith2021}. The Brain as a Query Interface (BQI) framework of Vu \cite{BQI} writes these dynamics as differential equations that link DMN activity, integrated information $\Phi$ and a reported quality of consciousness (QoC), fits their rates and gains to meta-analytic data, and proposes to steer them with a closed-loop brain--computer interface \cite[Sec.~7.2, 9.3 and 10.1]{BQI}; it is the theoretical foundation of the present study (Section~\ref{sec:theory}). Models of this kind promise quantities that can be estimated from data, such as the rate at which the DMN is suppressed and the weight of each latent state in the report, and they are put forward as the basis of model-based neurofeedback \cite{Sitaram2017}.

Whether these quantities can be estimated from the data that meditation studies actually collect, continuous EEG and a few self-report probes per session, is a separate question. In a companion paper \cite{companion} we replaced the BQI equations by a linear-Gaussian state-space model (M1) and used its exact expected Fisher information to compute how many subjects, how long a session, how frequent a probe and how clean an EEG feature are needed to estimate each parameter to a given precision. That design analysis was validated on simulated data only. No study has yet tested such a design analysis on real recordings, compared the estimates with the values that BQI derived from the literature, or asked whether the estimated parameters transfer from one laboratory to another (Section~\ref{sec:gap}).

These questions are hard to answer cleanly. The EEG features that stand in for the DMN and for $\Phi$ are proxies without a calibrated mapping to the state variables; a third of the recording can be missing; the model may be misspecified in ways that a bound cannot anticipate; and an analysis with many free choices (features, artefact rules, null models, exclusion criteria) can support almost any conclusion after the fact. We therefore preregistered the whole analysis plan for each of two open datasets, locked it with a SHA-256 hash before the EEG was examined \cite{Nosek2018}, organised the pipeline as a sequence of machine-checked gates that refuse to proceed when a prerequisite fails, and report every locked analysis whether or not it supported the predictions.

\subsection{Research Questions}
\begin{enumerate}
  \item[RQ1] \emph{Precision.} On a real experience-sampling design, with irregular probes, missing epochs and artefact rejection, is the relative error of the hardest parameter, the DMN rate constant $\kinh$, close to the Cram\'er--Rao bound (CRB) computed from the real probe schedules before the EEG is seen?
  \item[RQ2] \emph{Prediction.} Does a model that is identifiable in principle, fitted to continuous EEG and sparse self-reports, predict the reports out of sample better than models that ignore the dynamics?
  \item[RQ3] \emph{Agreement with BQI.} Do the estimated rate constants, coupling and gains agree with the values that BQI fitted to meta-analytic data?
  \item[RQ4] \emph{Replication.} Do the estimated parameters mean the same thing across recordings made in different laboratories with different amplifiers, montages, probe schedules and rating scales?
\end{enumerate}

\subsection{Contributions}
\begin{enumerate}
  \item The first test, to our knowledge, of a Fisher-information design analysis for a dynamical model of meditation against real experience-sampling EEG, on two independent open datasets, ds001787 \cite{Brandmeyer2018,Delorme2024ds} and OSF 3uszv \cite{RodriguezLarios2021}, under two locked preregistrations and eight verification gates.
  \item A reporting scheme of three error numbers, the CRB computed before the lock, the Fisher standard error at the estimate, and the parametric-bootstrap error, together with the observed-information error and a comparison with data simulated from the fit. On both datasets the first three agree for the DMN rate constant, and the fourth exposes a misspecification of the alpha-power channel \cite{White1982}.
  \item A preregistered out-of-sample test of the model against three null models, which shows that neither EEG proxy carries information about the reports beyond the reports' own serial dependence, on either dataset.
  \item A comparison of the estimates with the parameter values of BQI, and preregistered replication checks between the datasets, which show that the rate constants and the group difference between experienced and novice meditators do not transfer, together with a simulation that explains why: an Ornstein--Uhlenbeck fit returns the dominant time scale of the feature it is given.
  \item A fully reproducible pipeline (one command per step, checkpoints, 114 unit tests) with both preregistrations, locks and gate records.
\end{enumerate}

The rest of the paper follows the usual structure of a research article. Section~\ref{sec:litreview} reviews the literature on dynamical models of meditation, the BQI framework, EEG markers of DMN activity and of integration, the identifiability and design of dynamical models, and preregistration, and states the research gap. Section~\ref{sec:theory} sets out the theoretical framework: the BQI model with its equations, the derivation of the estimable model M1 and the predictions of BQI for the estimated quantities. Section~\ref{sec:methods} describes the methodology: study design, model, hypotheses, datasets, preregistration, preprocessing and statistical analysis. Section~\ref{sec:results} reports the results, Section~\ref{sec:discussion} discusses them together with the limitations of the study, and Section~\ref{sec:conclusion} concludes and outlines future work.
